\begin{abstract}
Penelitian ini membandingkan karakteristik bass, piano, gitar, brass,
dan tuba dalam aransemen ulang soundtrack \textit{Bad Piggies}.
Karena FL Studio tidak tersedia, data diperoleh dengan membaca not dan
susunan Playlist proyek FLP, mengekspornya ke MIDI, lalu merendernya
menggunakan FluidSynth dan SoundFont GeneralUser GS. Sebanyak 2219 not
berhasil direkonstruksi pada timeline yang sama. Objek penelitian adalah
komponen aransemen hasil rekonstruksi, bukan audio yang identik dengan
preset asli FL Studio. Analisis menggunakan representasi waktu, RMS,
DFT yang dihitung melalui FFT, serta STFT dengan jendela Hann 1024 dan
4096 sampel. Bass mempunyai sekitar $97{,}61\%$ energi di bawah 250 Hz,
sedangkan brass mempunyai sekitar $55{,}70\%$ energi di atas 1000 Hz.
Perbedaan waktu aktif menyebabkan RMS seluruh rekaman dan RMS saat
aktif memberikan gambaran berbeda. STFT memperlihatkan pergantian not
dan komponen yang bertahan, yang tidak langsung tampak pada spektrum
global. Variasi panjang jendela memperlihatkan kompromi lokalisasi
waktu dan frekuensi yang sesuai dengan Prinsip Ketaktentuan Heisenberg.
Konsistensi energi diperiksa dengan Parseval. Kanal DRUMS dan VOCALS
turut dianalisis dari WAV hasil dekode MP3. Konversi tersebut tidak
meningkatkan kualitas sumber; kedua kanal tetap menjadi pembanding
independen karena timeline berkas penuh belum diselaraskan dengan
rekonstruksi. Aktivitas VOCALS terkonsentrasi pada beberapa interval,
berbeda dari DRUMS yang aktif sepanjang sebagian besar rekaman.
\end{abstract}
\noindent\textbf{Kata kunci:} Fourier, FFT, STFT, Hann, instrumen musik,
MIDI, rekonstruksi audio.

\section{Pendahuluan}
Dua suara dapat memiliki amplitudo mirip, tetapi karakter berbeda.
Perbedaannya dapat berasal dari tinggi nada, susunan harmonik, serta
cara bunyi mulai dan meluruh. Sebagai analogi, bunyi rendah dapat
dibayangkan seperti ``dum'', sedangkan komponen tinggi dapat
mengingatkan pada desis atau ``cring''. Analogi ini hanya membantu
memahami frekuensi; bukan hasil pengamatan pendengaran terhadap setiap
berkas penelitian, dan tidak menyatakan bahwa frekuensi tertentu
menentukan suasana musik.

Objek penelitian berasal dari paket
\texttt{[ FREE FLP ] Bad Piggies - Theme NEW REMAKE FL STUDIO}.
Asal paket dirujukkan oleh penyedia data pada video Rick Ruyz
\cite{youtube}. Judul dan kanal diperiksa melalui metadata publik
YouTube. Paket memuat proyek FLP, dua MP3, dan sampel pendek. FLP
menyimpan informasi proyek, bukan rekaman audio siap analisis
\cite{flp}. Karena itu, diperlukan rekonstruksi sebelum komponen
instrumennya dapat dibandingkan.

Lima kelompok yang notnya terbaca adalah bass, piano, gitar, brass,
dan tuba. Label instrumen berasal dari channel sumber dan pemetaan
preset pengganti, bukan identifikasi otomatis. Penelitian tidak memakai
pembelajaran mesin dan tidak memisahkan instrumen dari campuran audio.
MIDI menyediakan not serta waktu kejadian, sedangkan SoundFont
menyediakan model bunyinya. Perbedaan dari FLEX/DirectWave merupakan
batas utama kesimpulan. Sebagai pembanding tambahan, kanal bernama
DRUMS dan VOCALS dibaca dari audio sumber yang sudah tersedia.
Keduanya tidak direkonstruksi melalui MIDI karena sumbernya berupa
klip audio, bukan not pada pattern.

FFT digunakan untuk merangkum kandungan frekuensi seluruh rekaman.
STFT menunjukkan kapan kandungan itu muncul atau berubah. Dua panjang
jendela dipakai agar pengaruh pemilihan potongan waktu dapat diamati
pada data yang sama. Penjelasan matematis mengikuti Kuliah 4--6 pada
folder \texttt{Materi}, termasuk sampling, Parseval, STFT dan
Heisenberg \cite{kuliah4,kuliah5,kuliah6}.

\section{Dasar Teori}
\subsection{Sinyal Diskret dan Notasi}
Audio digital menyimpan amplitudo pada waktu yang berjarak tetap.
Frekuensi sampling menyatakan banyak sampel per detik, berbeda dari
frekuensi nada dalam audio. Notasi pada Tabel~\ref{tab:notasi}
mengikuti bahan kuliah. Frekuensi Hz menggunakan $\nu$, sedangkan
$f$ menyatakan sinyal.
\begin{table}[H]
\centering\small
\caption{Simbol utama dan pengertiannya.}\label{tab:notasi}
\begin{tabular}{p{.22\linewidth}p{.7\linewidth}}
\toprule
Simbol & Pengertian \\
\midrule
$f(t)$, $f_k$ & Sinyal kontinu dan sampel ke-$k$ dari sinyal analisis. \\
$f_s$, $\Delta t$ & Frekuensi sampling dan jarak waktu sampel, $\Delta t=1/f_s$. \\
$N$, $T$ & Jumlah sampel dan durasi blok, $T=N\Delta t$. \\
$\nu$, $\omega$ & Frekuensi Hz dan frekuensi sudut rad/s; $\omega=2\pi\nu$. \\
$\widehat f$, $F_m$ & Transformasi Fourier kontinu dan koefisien DFT ke-$m$. \\
$w$, $M$, $H$ & Jendela, panjang jendela dan langkah pergeseran dalam sampel. \\
$\tau$, $q$, $Z_{q,m}$ & Pusat jendela, indeks bingkai, DFT potongan berjendela. \\
$V_wf$, $\Spec_wf$ & STFT dan spektrogram $|V_wf|^2$. \\
$P_m$, $P_{q,m}$ & Kerapatan daya satu sisi per Hz, global dan lokal. \\
$\sigma_t$, $\sigma_\omega$ & Penyebaran energi jendela terhadap waktu dan frekuensi sudut. \\
\bottomrule
\end{tabular}
\end{table}
Sampel berada pada $t_k=k\Delta t$, $k=0,\ldots,N-1$.
Dengan $f_s=44\,100$ Hz, batas Nyquist adalah $22\,050$ Hz.
Data diskret tidak menentukan secara unik komponen analog di atas
batas ini tanpa asumsi tambahan \cite{kuliah5}.

\subsection{Root Mean Square (RMS)}
RMS membantu membandingkan kekuatan amplitudo pada dua bagian audio.
Rerata amplitudo biasa kurang sesuai karena nilai positif dan negatif
dapat saling meniadakan. Untuk potongan berisi $Q$ sampel,
\begin{equation}
 \RMS=\sqrt{\frac1Q\sum_{k=0}^{Q-1}|f_k|^2}.
 \label{eq:rms}
\end{equation}
Kuadrat membuat kontribusi kedua tanda positif; rerata merangkum
potongan; akar mengembalikan satuan amplitudo. RMS tinggi berarti
amplitudo efektif lebih besar. Namun, RMS tidak langsung sama dengan
keras suara yang dirasakan manusia, sebab persepsi juga dipengaruhi
frekuensi, durasi dan kondisi pendengaran.

\subsection{Transformasi Fourier, DFT, FFT, dan Bin Frekuensi}
Fourier menyatakan sinyal melalui kontribusi osilasi pada berbagai
frekuensi. Materi memakai normalisasi kontinu uniter \cite{kuliah4}:
\begin{equation}
 \widehat f(\omega)=\frac1{\sqrt{2\pi}}\int_{\R}f(t)e^{-i\omega t}\,dt.
 \label{eq:ft}
\end{equation}
\textbf{DFT adalah transformasi matematis untuk data diskret, sedangkan
FFT adalah algoritma cepat untuk menghitung DFT}. Keduanya bukan dua
transformasi berbeda. DFT dan inversnya adalah
\begin{equation}
 F_m=\sum_{k=0}^{N-1}f_ke^{-2\pi i mk/N},\qquad
 f_k=\frac1N\sum_{m=0}^{N-1}F_me^{2\pi i mk/N}.
 \label{eq:dft}
\end{equation}
Satu \textit{bin} adalah satu titik frekuensi tempat DFT dihitung.
Pada sisi nonnegatif, $\nu_m=mf_s/N$ dan
\begin{equation}
 \Delta\nu=\frac{f_s}{N},\quad \omega_m=2\pi\nu_m,\quad
 \widehat f(\omega_m)\approx\frac{\Delta t}{\sqrt{2\pi}}F_m.
 \label{eq:grid}
\end{equation}
Jarak bin kecil berarti titik pengamatan rapat. Kemampuan memisahkan
dua komponen tetap dipengaruhi panjang pengamatan, jendela dan sifat
sinyal. Faktor terakhir menghubungkan keluaran FFT tanpa normalisasi
maju dengan konvensi kontinu materi \cite{kuliah5,numpy}.

Grafik memakai kerapatan daya satu sisi
\begin{equation}
 P_m=\frac{c_m|F_m|^2}{f_sN},
 \label{eq:psd}
\end{equation}
dengan $c_m=2$ pada bin positif yang mempunyai pasangan negatif,
dan $c_m=1$ pada DC serta Nyquist jika ada. Daya pita dihitung dengan
menjumlahkan $P_m\Delta\nu$. Tinggi satu puncak berbeda makna dari
jumlah energi dalam sebuah rentang frekuensi.

\subsection{Identitas Parseval}
Jika perubahan representasi dilakukan konsisten, energi dari sampel
dan koefisien Fourier harus sama. Plancherel memberikan
$\int|f(t)|^2dt=\int|\widehat f(\omega)|^2d\omega$ \cite{kuliah4}.
Versi diskret untuk memeriksa komputasi ialah
\begin{equation}
 E_t=\Delta t\sum_k|f_k|^2
 =\frac{\Delta t}{N}\sum_{m=0}^{N-1}|F_m|^2=E_F.
 \label{eq:parseval}
\end{equation}
Pada FFT satu sisi digunakan bobot $c_m$. Energi bersatuan amplitudo
digital kuadrat dikali detik, bukan joule, karena audio tidak
dikalibrasi sebagai tekanan atau daya akustik.

\subsection{Short-Time Fourier Transform (STFT) dan Spektrogram}
FFT global membaca seluruh audio sekaligus. STFT menggeser potongan
pendek, lalu menghitung Fourier berulang kali. Frekuensi kemudian
dihubungkan dengan lokasi waktunya. Definisi materi adalah
\cite{kuliah6}
\begin{equation}
 V_wf(\tau,\omega)=\frac1{\sqrt{2\pi}}
 \int_{\R}f(t)\overline{w(t-\tau)}e^{-i\omega t}\,dt,
 \qquad \Spec_wf=|V_wf|^2.
 \label{eq:stft}
\end{equation}
Sumbu horizontal spektrogram menyatakan waktu, vertikal frekuensi,
dan warna kekuatan komponen menurut skala daya. Garis mendatar dapat
menunjukkan frekuensi yang bertahan; guratan tegak dapat menunjukkan
kejadian singkat yang menjangkau banyak frekuensi. Pola tersebut
perlu dikaitkan dengan data, bukan langsung dianggap identitas instrumen.

Pemotongan tajam pada tepi dapat menimbulkan kebocoran spektral:
kontribusi komponen menyebar ke frekuensi sekitar. Hann mengurangi
bobot tepi. Varian periodik yang dipakai adalah
\begin{equation}
 w_r=\frac12-\frac12\cos\left(\frac{2\pi r}{M}\right),
 \quad r=0,\ldots,M-1.
 \label{eq:hann}
\end{equation}
Kedua panjang jendela memakai bentuk sama agar perbandingan adil
\cite{hann}. Hann mengurangi kebocoran, tetapi tidak menghilangkannya
dan mempunyai respons frekuensi dengan lebar tertentu.

\subsection{Prinsip Ketaktentuan Heisenberg}
Untuk mengetahui kapan bunyi berubah dengan lebih tepat, diperlukan
potongan pendek. Namun, potongan pendek memuat lebih sedikit siklus
sehingga frekuensi berdekatan lebih sulit dibedakan. Potongan panjang
memperinci frekuensi, tetapi merangkum rentang waktu lebih luas.

$\sigma_t$ mengukur penyebaran energi jendela dari pusat waktunya,
sedangkan $\sigma_\omega$ mengukur penyebaran energi transformasi
jendela dari pusat frekuensinya. Materi memberikan
\begin{equation}
 \sigma_t\sigma_\omega\geq\frac12,\qquad
 \sigma_t\sigma_\nu\geq\frac1{4\pi}.
 \label{eq:heisenberg}
\end{equation}
Batas pertama memakai rad/s, yang kedua Hz. Durasi $M/f_s$, jarak
bin $f_s/M$ dan langkah $H/f_s$ bukan simpangan baku tersebut.
Definisi integral dan perhitungan Hann diberikan pada
Lampiran~\ref{app:normalisasi}.

\section{Permasalahan}
Nada rendah yang dipetik dan ditiup dapat memiliki frekuensi dasar
berdekatan, tetapi susunan harmonik dan perubahan amplitudo berbeda.
Sebaliknya, instrumen sejenis dapat berbeda spektrum karena memainkan
register atau not berbeda. Penelitian menanyakan apa yang dapat
dijelaskan dari komponen aransemen yang tersedia.
\begin{enumerate}
\item Bagaimana perbedaan waktu aktif, amplitudo efektif, awal dan
peluruhan bunyi kelima komponen rekonstruksi?
\item Frekuensi mana yang berkontribusi kuat dan bagaimana energi
terdistribusi dalam pita frekuensi masing-masing komponen?
\item Informasi waktu apa yang ditambahkan STFT terhadap FFT global?
\item Bagaimana Hann 1024 dan 4096 sampel mengubah tampilan kejadian
singkat dan komponen frekuensi berdekatan?
\item Bagaimana hasil sesuai dengan Fourier, Parseval dan Heisenberg,
serta mendukung deskripsi nonstasioner?
\end{enumerate}
Perbedaan yang ditemukan berlaku untuk kombinasi partitur, velocity,
register dan preset ini. Analisis tambahan membandingkan karakter
waktu dan frekuensi kanal DRUMS dan VOCALS. Kedua kanal merupakan
pembanding terpisah; timecode berkas penuhnya belum disetarakan dengan
lima WAV rekonstruksi.
